% ============================================================
% SECTION: Automated Lab Provisioning with Cloud-Init
% ============================================================

\section{Automated Lab Provisioning}
\label{sec:cloud-init}
% ------------------------------------------------------------

A key reproducibility requirement for the course is that each student works in an
identical, isolated EVE-NG environment. Manually cloning and reconfiguring virtual
machines for each student is error-prone and does not scale. To address this,
a \textit{cloud-init}-based provisioning system was designed and verified on top
of Proxmox~\cite{proxmox}, enabling an instructor to deploy a ready-to-use
EVE-NG instance for each student with a single command.

\subsection{Architecture}

The provisioning model follows the standard Proxmox linked-clone pattern with
a cloud-init drive:

\begin{itemize}
    \item \textbf{VM~113 (template)}: A prepared EVE-NG base installation kept in
          a stopped state. Its disk is stored as a ZFS zvol on the Proxmox pool,
          which allows efficient copy-on-write linked clones without duplicating the
          full 200~GB image.
    \item \textbf{Snippets directory} (\path{/var/lib/vz/snippets/}): Contains one
          \texttt{cloud-config} YAML file per student. Each snippet sets the
          hostname, creates the student user with a dedicated SSH key, and configures
          the \texttt{pnet0} bridge address via a \texttt{runcmd} block.
    \item \textbf{Student VM (linked clone)}: Created on demand from the template.
          Proxmox generates a small ISO (\textit{cidata}) at boot from the assigned
          snippet and mounts it at \path{/dev/sr0}, where cloud-init reads it.
\end{itemize}

\subsection{Cloud-Init Configuration on the Template}

Three configuration files are placed on VM~113 before it is sealed as a template:

\begin{itemize}
    \item \textbf{\path{/etc/cloud/ds-identify.cfg}}: Sets
          \texttt{policy: enabled,found=all} so that cloud-init does not
          self-disable when no datasource is found on the first second of boot ---
          a common failure mode in non-cloud images.
    \item \textbf{\path{99\_proxmox.cfg}}: Configures
          \texttt{datasource\_list: [NoCloud, ConfigDrive, None]} and sets the
          \textit{fs\_label} to \texttt{cidata}, which matches the label Proxmox
          stamps on the generated ISO.
    \item \textbf{\path{99\_disable\_network.cfg}}: Sets \texttt{network: config:
          disabled}. This is critical: EVE-NG manages its own \texttt{pnet0--pnet9}
          bridge stack, and allowing cloud-init to reconfigure the network interface
          would destroy those bridges on every boot.
\end{itemize}

\subsection{Student Snippet Structure}

Each per-student YAML snippet follows this structure:

\begin{lstlisting}[caption={Cloud-init user-data snippet for a student VM}, label={lst:cloud-init-yaml}]
#cloud-config
hostname: eve-ng-alumno2
manage_etc_hosts: true

users:
  - name: alumno2
    shell: /bin/bash
    sudo: ALL=(ALL) NOPASSWD:ALL
    ssh_authorized_keys:
      - ssh-ed25519 AAAA...key... alumno2@eveng-lab

chpasswd:
  list: |
    alumno2:Lab2024!
  expire: false

runcmd:
  - /bin/bash -c 'ip addr flush dev pnet0;
    ip addr add 192.168.0.211/24 dev pnet0;
    ip route add default via 192.168.0.1 dev pnet0;
    sed -i "s/address 192.168.0.[0-9]*/address 192.168.0.211/"
    /etc/network/interfaces'

package_upgrade: false
\end{lstlisting}

The \texttt{runcmd} block flushes and re-assigns the \texttt{pnet0} bridge IP
to the student-specific address and persists it in
\path{/etc/network/interfaces} for subsequent reboots. Using a dedicated SSH
key per student ensures that no credentials are shared between instances.


\newpage
\subsection{Provisioning Workflow}

Once the template and snippets are in place, deploying a new student instance
requires four commands on the Proxmox host:

\begin{lstlisting}[language=bash, caption={Provisioning a student EVE-NG VM from the template}]
# 1. Linked clone from template (no full disk copy)
qm clone 113 <VM_ID> --name eve-ng-alumnoN --full 0

# 2. Assign the student snippet and IP
qm set <VM_ID> --cicustom \
    'user=local:snippets/userdata_alumnoN.yaml'
qm set <VM_ID> --ipconfig0 'ip=192.168.0.2XX/24,gw=192.168.0.1'
qm cloudinit update <VM_ID>

# 3. Start the VM
qm start <VM_ID>

# 4. Connect after ~40 seconds
ssh alumnoN@192.168.0.2XX -i /path/to/alumnoN_key
\end{lstlisting}

The linked-clone model means only the delta from the base image is stored per
student, keeping storage overhead proportional to actual usage rather than
multiplying the full image size.

\subsection{Verification}

The complete provisioning pipeline was validated with a fresh VM~114 deployment
from the sealed template. A checklist of 28 verification points was used, covering:
template integrity, cloud-init configuration, clone creation, network assignment,
user creation, SSH access, and EVE-NG service health. All 28 checks passed
with a VM reaching full operational state --- including a working EVE-NG web
interface --- within approximately 35 seconds of \texttt{qm start}.

\subsection{Problems Encountered}

\subsubsection{LVM Filter Blocking ZFS Zvol Mount}

Mounting the template disk from the Proxmox host for inspection required
activating the LVM volume group on the ZFS zvol. Proxmox's default
\path{lvm.conf} applies a \texttt{global\_filter} that excludes
\path{/dev/zd.*} devices to prevent conflicts. The workaround was to activate
the specific volume group using \texttt{vgchange --select 'uuid=<UUID>'}
instead of modifying the global filter, preserving the host's LVM configuration
while allowing temporary access to the template disk.

\subsubsection{IP Persistence Across Reboots}

The initial runcmd block updated the in-memory bridge address but did not persist
the change. On the second boot, the VM reverted to the template's base IP
(\texttt{192.168.0.133}). The fix was to include a \texttt{sed} command in the
same runcmd block that replaces the address in
\path{/etc/network/interfaces} atomically with the student IP, so the correct
address is applied by \texttt{ifup} on every subsequent boot without requiring
cloud-init to run again.

\subsection{Further Reference}

Full step-by-step instructions, complete configuration file listings, and the
28-point verification checklist are available in the online lab documentation:
\url{https://theegea.github.io/TFG/deployment/cloud-init/}.
